\chapter{Classification}

\section{Sample Data}

The training and testing data come from a publicly available data set containing
customer reviews of books purchased in Amazon e-shops. That this data can be
freely downloaded from a publicly accessible URL \cite{ref190}. Experiments, demonstrated below, use a randomly selected set from the large original database. The
data collection for demonstrating capabilities of selected classifiers are limited
here to just 1000 samples which are evenly split into two classes: books with
positive reviews (five-star reviews) and books with negative reviews (one-star
reviews). All reviews are written by customers in English.

To illustrate the original reviews, here are some randomly selected data items
from the data set coming from the URL mentioned above \cite{ref190} and used in this
book as input for the presented classification algorithms. These demonstration
reviews are not related to any particular book. To save space, long reviews have
not been included here.

First, five positive reviews:

\begin{itemize}
    \item Really interesting. Everything you need, the most important is inside. Clear explanation, good pictures to represent ’the move’. This books helps me to improve
myself and I will soon read it again... I guess it can be useful for beginner too...
    \item This book is good on many levels. I learned about Ethiopia, it’s culture and the
struggles during this time. Excellent writing.

    \item Ron Hansen is among the best writers delving into unusual real-life historical situations, and trying to make sense of them. Mariette is a pretty 17 year old drawn
to the convent, who immediately becomes a special, vision influenced, faith inspired, stigmatized True Believer. Many of the older nuns are skeptical (to say the least), while the younger ones are in awe. Mr. Hansen guides us right into the
heart and soul of this Catholic community, and of course Mariette. After reading this book, one feels a certain understanding of this potential saint, and how
religion may become an obsession. I would personally agree with her physician
/ father, the scientific skeptic, but Mariette’s story is certainly believable in the
hands of a fine writer like Mr. Hansen!

    \item Excellent book by a wonderful leader. Used it to develop a presentation on leadership for a mentoring team. Practical and excellent lessons for everyone. Thank
you, Mr. Powell.

    \item Harry Stein writes one of the most poltically honest memoirs I’ve read in a long
time. In straightforward speak that skirts no liberl ideals, Stein chops down an
ideology that he swore by in the 1960s but now sees as a hinderance to our
modern-day culture. While his liberal foes cringe, every independently-minded
American should give Stein’s views a chance. For the naysers of this book, especially hardcore liberals, I found it interesting that many of the facts Stein used to
back up his positions were pulled out of the liberally-slanted media he so accurately portrays.
\end{itemize}

Next, five negative reviews:

\begin{itemize}
    \item For people who have some knowledge of China, the book is of NO interest. Other
books by James Mann (About Face), Jonathan Fenby, not to forget Spence (In
Search of Modern China), are clearly superior.Kissinger’s admiration for Mao
is a disgrace. Mao was a murderer and Kissinger’s repetitions admiration for
China’s civiliztion is non-sence. Over the centuries, China has been a mess, with
short periods of success, but mostly calamities. Because of its size, China should
have been a success, but it is not. Japan was must faster on its feet! China’s
isolation was a huge mistake. It is only recently, thanks to Deng and Western
support, that China at last is emerging!Kissinger acknowledges that China has
always played one ”barbarian” against the other. It is very true but he does not
draw any conclusion. To-day China is gaming both Europe and America: Just
look at the ”free” transfer of thechnology. The West is stupid and we should
unite ton face China...The last point is Climate Change. To write about China and
not discuss the catastrophic impact of climate change, is a major failure of the
book. China - and Asia - will be destroyed by climate change.The combination
of bad governance, corruption, and climate change is DEADLY. It is a Chinese
and Asian disease.

    \item Let me set the scene for you:I read the first three chapters of this book a few
years ago right before it was published, and because I was immediately hooked
and knew that I’d blow through the book in a day (then end up miserable until
the next one was published over a year later), I decided to put it aside until there
were more books to read. So fast forward to last week... I had an eight hour shift taking inventory at a local retailer scheduled, so I pulled up my Kindle
wishlist and looked for something that was also available on Audible. This book
was the winner, and–shame on me–I was, unfortunately, too excited to bother
with listening to the Audible sample first. Oops. Big OOPS.Within thirty seconds
I was asking myself if the producers of the audiobook were serious. By fortyfive seconds, I was wondering if the author had approved the narrator and was
happy with her. Then, around the sixty second mark, I decided it didn’t matter
to me. I turned it off. I couldn’t listen to a story being read that way, no matter
how compelling it is. So here I am; warning others away. READ the book–don’t
bother trying to listen to it. Or, if listening is your only option, then I would
suggest looking into using Ivona Voices to create an MP3 version of the Kindle
book. You would be better off, in my not so humble opinion .

    \item Foolishly, I purchased a copy of this book for a young mother, at her request.Now,
I am familiar with the contents of said book, and I don’t know how to deal with
the guilt. Oh, my God! What can I do to protect the daughter of this young
woman? If she is naive enough to swallow this ”author’s” militant sickness, I
have doomed an innocent child to a life of horror.On a more positive note, if I
can convince the recipient to toss the volume into the dumpster, I will have removed one copy of it from circulation, and perhaps saved another innocent child
from this abuse.This abomination is not recommended.

    \item The Lombardi Rules attempts to provide a philosophical background of Coach
Lombardi’s beliefs and apply them to life and business. This was written by Vince
Lombardi Jr, not the man himself, and although you’d expect his son to also possess many of the inspiring attributes of Coach Lombardi, they are not exhibited
in this book. It is essentially a collection of Vince Lombardi quotes turned into
2-3 page chapters. Making matters worse is that the text supporting the quotes
ranges from cliche to painfully obvious with nothing profound or meaningful,
not even a decent example from Lombardi’s illustrious coaching career. Its always frustrating when you could read the chapter headings and gain as much
value as the entire book. You would be better served just searching for Vince
Lombardi quotes than wasting any time with this book.

\item I normally adore anything and everything Victoria Alexander. This book started
out to be good but I felt it quickly left me feeling bored and dissatisfied. The
Princess’s character was very annoying. Although I love Weston I just couldn’t
get passed how slow the book moved along. I would not recommend this book
to a first time reader of Victoria Alexander. She is a wonderful writer, this book I
feel was just a lemon out of hundreds.
\end{itemize}

Half of the examples used for the demonstration of classifiers, 500, are labeled as positive reviews, the remaining half as negative. All reviews have
been taken as they are, including possible grammar mistakes, typing errors, and so on, as typical for contributions on social networks. No special adjustment was done except for removing irrelevant attributes (those that do not contribute to classification). Removing such attributes decreases computational complexity (by lowering the number of dimensions, that is words, therefore saving memory and time). In the package \texttt{tm}, there is a very efficient function,
removeSparseTerms(), reducing matrix sparsity (removing columns containing too many zeros according to parameters given by a user). The parameter was
experimentally set to 0.995. This means that attributes that appear in at least 0.5\%
(1-0.995) of documents are retained. This setting filtered a majority of attributes
(only 2218 left, out of 17903) while retaining the classification accuracy.

The data are prepared in the pure textual form and their input format is CSV
(comma separated values), including the first line which contains names of attributes (which are the words themselves). As for the format of the input data, it
is represented simply as a vector for each review, where the values of elements
are the word frequencies in individual reviews.

In a simplified way, vectors can be viewed like genomes composed of genes,
where a genome constitutes a textual document and genes represent words. The
more genes are shared between documents, the more similar (‘related’) the documents are. A vector direction points at a subspace containing similar documents.

The data representation is a matrix – every column represents an attribute,
every row (vector) represents a textual review. The last vector element is a class
label (which is a discrete value). In each of the following sections, the same
demonstration data set is used to present all selected classifiers mentioned.

\section{Selected Algorithms}

In this book, a chosen set of algorithms consists of typical ones employed for
text mining up to the present. Naturally, the machine-learning selected here algorithms do not cover every possibility and a reader is encouraged to experiment
with many other existing methods. One of the problems of machine-learning is
that there seems to be no universal, best, algorithm for any data type. In practice,
it is often necessary to experimentally look for the best data/text mining tool.

This problem has been discussed in many publications; for example, refer to
\cite{ref281} and \cite{ref282}. Finding and presenting the final optimal algorithm as a solution
for a given application might sometimes just be that proverbial tip of the glacier,
while the rest can absorb the overwhelming majority of invested time and other
resources.
This part of the book presents machine-learning algorithms which employ
various principles of learning:

\begin{itemize}
    \item Optimal and Na¨ıve Bayes’ Classifier – traditional and fast methods based
on (conditional) probability theory, with excellent mathematical support,
including uncertainty;
    \item k-Nearest Neighbors – methods based on similarity with known entities,
very fast during the training phase, might be appreciably slow during the
classification process;
    \item Decision Trees – popular ‘divide and conquer’ algorithms, which can
provide understandable knowledge which may be transferred from a tree
to a set of rules;
    \item Random Forest – classification by a group of decision trees, which usually provides good results based on a noticeable reduction of possible tree
correlations and over-fitting;

    \item Support Vector Machines (SVM) – an often applied algorithm to problems where it is convenient to look for linear boundaries between classes
in artificially enlarged spaces when the original space has just non-linear
boundaries; SVM can give very good results but might have a high computational complexity;

    \item Adaboost – a favorite method based on using a set of ‘weak learners’
that, together, can provide excellent results, even if each of the learners
classifies with an error only slightly less than 50%.

    \item Deep Learning – one of very modern methods employing the well-known
and time-proven artificial neural networks in many applications. This approach relies on the power of multilayered networks, where the deepness
is given by the number of those layers; each layer makes new combinations of input values, thus creating an abstract space where the classification can work well. Deep Learning is a resource (time, memory)
demanding method, which can, today, take advantage of powerful hardware.
\end{itemize}

All the above mentioned algorithms have many implementations, including
the popular language R. Sometimes, in the same programming language or a system, a user can select more or less different implementations according her or his
application needs. One of the most important matters is to learn which parameters are available for an algorithm and which values it accepts, including the
default values. The default algorithm parameters are often set to the statistically
most used values; however, it does not have to be the optimal ones.

In this book, a reader cannot find a detailed description of all possible parameters because the intention was not to supply a manual. However, the detailed
understandable descriptions, including samples of R-code, are available from the Internet as PDF documents and may be easily downloaded from the following
web pages:

\begin{itemize}
    \item Optimal and Na¨ıve Bayes’ Classifier – see \cite{ref182},
    \item k-Nearest Neighbors – see \cite{ref226},
    \item Decision Trees – see \cite{ref154},
    \item Random Trees and Forest – see \cite{ref171},
    \item Support Vector Machines (SVM) – see [153] and \cite{ref193},
    \item Adaboost – see \cite{ref60},
    \item Deep Learning – see \cite{ref83}.
\end{itemize}

\section{Classifier Quality Measurement}

After selecting and training a classifier algorithm, the supervisor needs to know
the classifier efficiency – based on the training samples, is the mined knowledge
sufficient enough for its successful application to future data?

Once the classifier training process has reached a certain peak value (depleted
time, or maximum number of steps, or any other appropriate criterion), which
completes the learning, it is necessary to express the algorithm’s training quality, preferably numerically. This value can also be used to compare the mutual
performance degree of various classifiers under investigation over the same data
from a given application point of view.

Several methods are used in practice, of which the four most commonly used
will be mentioned here. To evaluate the quality of a classifier, it is necessary to
use so-called test data (samples), which correspond thematically to the training
data, but were not used during training. It is similar to teaching at a school where
a teacher (supervisor) teaches pupils to correctly compute some examples, and
uses other, but similar examples to test their understanding.

From a binary perspective, four terms are used for the computation of classifier efficiency: 
\textit{TP} (True Positive) for the number of data items classified correctly
from a certain class \textit{C} point of view, 
\textit{FP} (False Positive) for items incorrectly assigned to \textit{C}, 
\textit{TN} (True Negative) for items correctly classified as not belonging to
\textit{C}, and 
\textit{FN} (False Negative) for items belonging to \textit{C} but assigned to other class
or classes~\cite{ref107}:

\begin{itemize}
    \item \textit{Accuracy}: 
$$
\textit{Accuracy} = ({\textit{TP} + \textit{TN}})/({\textit{TP} + \textit{TN} + \textit{FP} + \textit{FN}})
$$
that is, the number of correct answers to the number of all classified samples, often expressed as a percentage.
    \item \textit{Precision}: 
$$
\textit{Precision} = {\textit{TP}}/({\textit{TP} + \textit{FP}})
$$
that is, the number of items correctly labeled as belonging to the positive class \textit{C} divided by the total number of items labeled both correctly and incorrectly as belonging to \textit{C} (which is actually the number of relevant items in a set of items assigned to \textit{C} by the investigated classifier).
    \item \textit{Recall}: 
$$
\textit{Recall} = {\textit{TP}}/({\textit{TP} + \textit{FN}})
$$
that is, the proportion of actual positives that was identified correctly.
    \item \textit{F-measure}: it is the weighted harmonic mean of the precision and recall:
$$
F = \frac{1}{\alpha \frac{1}{\textit{precision}} + (1-\alpha) \frac{1}{\textit{recall}}}
= \frac{(\beta^2 + 1) \cdot \textit{precision} \cdot \textit{recall}}{\beta^2 \cdot \textit{precision} + \textit{recall}}
$$

where 
$
\beta^2 = \frac{1-\alpha}{\alpha}.
$
When $\beta$ increases, the precision has a higher importance. When the precision and recall are equally weighted (have the same importance), $\beta = 1$, and the measure is also called \textit{F1} or \textit{F-score}, 
which is the harmonic mean of the precision and recall:
$$
F1 = \frac{2 \cdot \textit{precision} \cdot \textit{recall}}{ 2 \cdot\textit{precision} + \textit{recall} }
= \frac{2 \cdot \textit{TP}}{2 \cdot \textit{TP} + \textit{FP} + \textit{FN}}.
$$

It is possible to draw a graph where the horizontal axis represents the
recall and the vertical axis represents the precision, however, sometimes
not a curve but just a number, which expresses the classification quality,
is needed.

\end{itemize}


Various metrics evaluating the results of classification can be calculated for
each class. For example, the \textit{precision} for class 0 is 0.5 and the \textit{precision} for
class 1 is 0.9. Often, we want just one number evaluating a classifier. Thus,
the next step might be to average these two values. Using \textit{macroaveraging}, which
simply relies on the arithmetic mean, the average precision is 0.7. Both values representing a classifier’s performance for each class have an equal weight.

\textit{Microaveraging} aggregates the decisions across all classes, taking every single classified item into consideration. Thus, the classes have different weights
depending on their size (bigger classes contribute to a higher measure value). When class 0 is predicted for 20 items and class 1 for 100 items (class 1 is predicted five times more frequently than class 0), microaveraged precision will be
$
{20}/({20+100})) \cdot 0.5 + {100}/({20+100}) \cdot 0.9 = 0.83.
$
It is obvious that the precision is then closer to the value of the precision for class 1, which is bigger.




